\section{Document Provenance and Clustered Significance Testing}
\label{app:provenance}
\label{app:clustering}

\subsection{Document Provenance}

The 192 documents listed in Appendix~\ref{app:sourcedocs} are the full collected corpus. The 406
annotated items trace to \textbf{34} of them, spanning both regulators; Table~\ref{tab:provenance}
summarises the distribution.

\begin{table}[h]
\centering\small
\begin{tabular}{lr}
\toprule
Documents collected (full corpus) & 192 \\
Documents represented in the 406 items & 34 \\
Largest single document (item count) & 63 \\
Median items per document & 6 \\
\bottomrule
\end{tabular}
\caption{Document provenance summary. Full item-level provenance
(\texttt{evaluation/provenance\_audit.csv}) is released with the benchmark.}
\label{tab:provenance}
\end{table}

\subsection{Document-Clustered Bootstrap}

Items sharing a source document are not independent draws. We construct resampling clusters as the
connected components of a graph over the 34 documents, with an edge added between
\texttt{document\_a} and \texttt{document\_b} for every contradiction-detection item---since a CON
item depends on two documents jointly, keeping them in separate clusters would let dependence leak
across the resampling unit boundary. This produced \textbf{27} connected components (sizes: 68, 48,
46, 41, 41, 35, 20, 13, 12, 11, 11, 10, 9, 7, 6, 4, 4, 4, 3, 3, 2, 2, 2, 1, 1, 1, 1); four are singleton
documents that appear in no contradiction pair.

We re-ran paired-bootstrap significance testing (10,000 resamples) at the document-cluster level for
\textbf{all 66 pairwise comparisons}, plus the human-reference comparison separately: resampling
clusters with replacement, and including every item in each drawn cluster's data, rather than
resampling items directly. Table~\ref{tab:clustered} reports the seven comparisons discussed in the
main text; the full 66-pair table is released with the evaluation code
(\texttt{evaluation/clustered\_bootstrap\_full66.json}).

\begin{table}[h]
\centering\small
\begin{tabular}{lrr}
\toprule
Comparison & Item-level $p$ & Cluster-level $p$ \\
\midrule
Gemini 2.5 Flash vs.\ Qwen3-32B & 0.052 & 0.060 \\
Gemini 2.5 Flash vs.\ Gemma 4 E4B & $<$0.001 & $<$0.001 \\
Llama 4 Scout 17B vs.\ LLaMA-3.3-70B & 0.788 & 0.825 \\
GPT-OSS 120B vs.\ GPT-OSS 20B & 0.904 & 0.936 \\
Gemma 4 E4B vs.\ DeepSeek-R1-Distill & 0.186 & 0.424 \\
Gemma 4 E4B vs.\ Mistral-7B & 0.254 & 0.431 \\
Human vs.\ Gemini 2.5 Flash (60-item subset, 17 clusters) & 0.427 & 0.558 \\
\bottomrule
\end{tabular}
\caption{Item-level vs.\ document-clustered bootstrap $p$-values for the seven comparisons discussed
in the main text. All seven agree in significance direction. Across the full 66-pair table, 29 of 66
are significant under clustering versus 36 item-level, and 9 pairs (not shown here, none discussed
in the main text) flip significance direction --- see the released full comparison for the complete
picture rather than only these seven.}
\label{tab:clustered}
\end{table}

Cluster assignment was checked for correctness before use: every item maps to exactly one component,
every document appears in exactly one component, and the procedure is deterministic under a fixed
random seed (verified by re-running with an identical seed and confirming bitwise-identical output).
We do not treat "the clustered interval is wider" as a correctness check---no such guarantee exists
in general--- and instead verified the construction directly.
